% ============================================================
%  Openbook / Formelsammlung Template — strikt nach my_dhbw
%  Stil: 2-spaltig, kompakt, scriptsize Tabellen, dünne Regeln.
%  Renderer: lualatex (zweimal, wg. Unicode-Inhalt)
% ============================================================
\documentclass[10pt, a4paper]{article}

\usepackage[ngerman]{babel}
\usepackage{fontspec}
\setmainfont[
  Path=/usr/share/fonts/TTF/,
  Extension=.ttf,
  BoldFont=DejaVuSans-Bold,
  ItalicFont=DejaVuSans-Oblique,
  BoldItalicFont=DejaVuSans-BoldOblique
]{DejaVuSans}
\setsansfont[
  Path=/usr/share/fonts/TTF/,
  Extension=.ttf,
  BoldFont=DejaVuSans-Bold,
  ItalicFont=DejaVuSans-Oblique,
  BoldItalicFont=DejaVuSans-BoldOblique
]{DejaVuSans}
\setmonofont[Path=/usr/share/fonts/TTF/,Extension=.ttf]{DejaVuSansMono}

\usepackage[top=0.8cm, bottom=0.8cm, left=0.8cm, right=0.8cm, footskip=0.4cm]{geometry}
\usepackage{amsmath, amssymb}
\usepackage{multicol}
\usepackage{xcolor}
\usepackage{titlesec}
\usepackage{enumitem}
\usepackage{fancyhdr}
\usepackage{lastpage}
\usepackage{microtype}
\usepackage{etoolbox}

\usepackage{booktabs}
\usepackage{array}
\usepackage{tabularx}
% ltablex/longtable bewusst NICHT geladen — kollidiert mit multicols.
\usepackage{colortbl}
\usepackage[most]{tcolorbox}
\usepackage{listings}
\usepackage[normalem]{ulem}
\usepackage{adjustbox}
\usepackage[colorlinks=false]{hyperref}

\renewcommand{\familydefault}{\sfdefault}

% ── Theme (my_dhbw accent) ────────────────────────────────────
\definecolor{accent}{RGB}{0, 60, 113}
\definecolor{primaryblue}{RGB}{0, 60, 113}
\definecolor{codebg}{RGB}{245, 245, 245}
\definecolor{figurebg}{RGB}{232, 241, 255}
\definecolor{tablehintbg}{RGB}{240, 255, 240}
\definecolor{solutionbg}{RGB}{242, 255, 242}
\definecolor{rulegray}{RGB}{180, 180, 180}
\definecolor{tableheadbg}{RGB}{0, 60, 113}

% ── Header/Footer — wie my_dhbw formelsammlung.tex ────────────
\pagestyle{fancy}
\fancyhf{}
\renewcommand{\headrulewidth}{0pt}
\renewcommand{\footrulewidth}{0.3pt}
\fancyfoot[L]{\tiny\color{accent}\nichttitle}
\fancyfoot[R]{\tiny\color{accent}\thepage\,/\,\pageref{LastPage}}

% ── Typografie — wie formelsammlung.tex ───────────────────────
\titleformat{\section}
    {\color{accent}\footnotesize\bfseries}{}{0pt}{}
    [\vspace{-2pt}\color{accent!50}\rule{\linewidth}{0.4pt}]
\titleformat{\subsection}
    {\scriptsize\bfseries\color{accent!85!black}}{}{0pt}{}{}
\titleformat{\subsubsection}
    {\scriptsize\itshape\color{accent!60!black}}{}{0pt}{}{}
\titlespacing*{\section}{0pt}{4pt}{1pt}
\titlespacing*{\subsection}{0pt}{3pt}{0pt}
\titlespacing*{\subsubsection}{0pt}{2pt}{0pt}

\setlength{\parindent}{0pt}
\setlength{\parskip}{1pt}
\renewcommand{\arraystretch}{1.0}
\setlist{noitemsep, topsep=0pt, parsep=0pt, leftmargin=1.1em}

\setlength{\columnsep}{6mm}
\setlength{\columnseprule}{0.2pt}

% Globaler scriptsize-Body. Tabellen zusätzlich auf scriptsize zwingen.
\AtBeginEnvironment{tabularx}{\scriptsize}
\AtBeginEnvironment{tabular}{\scriptsize}

% ── Boxen (kompakt) ───────────────────────────────────────────
\newtcolorbox{figurebox}[1]{
  colback=figurebg, colframe=accent!50,
  title={\scriptsize Abbildung: #1}, fonttitle=\scriptsize\bfseries,
  fontupper=\scriptsize, sharp corners,
  boxsep=2pt, left=3pt, right=3pt, top=2pt, bottom=2pt,
  before skip=3pt, after skip=3pt
}
\newtcolorbox{tablehintbox}[1]{
  colback=tablehintbg, colframe=green!50!black!50,
  title={\scriptsize Tabelle: #1}, fonttitle=\scriptsize\bfseries,
  fontupper=\scriptsize, sharp corners,
  boxsep=2pt, left=3pt, right=3pt, top=2pt, bottom=2pt,
  before skip=3pt, after skip=3pt
}
\newtcolorbox{solutionbox}[1]{
  colback=solutionbg, colframe=green!60!black!60,
  title={\scriptsize #1}, fonttitle=\scriptsize\bfseries,
  fontupper=\scriptsize, sharp corners,
  boxsep=2pt, left=3pt, right=3pt, top=2pt, bottom=2pt,
  before skip=3pt, after skip=3pt
}

\lstset{
  basicstyle=\ttfamily\tiny,
  backgroundcolor=\color{codebg},
  breaklines=true,
  frame=none,
  xleftmargin=2pt, xrightmargin=2pt,
  aboveskip=2pt, belowskip=2pt,
  keepspaces=true
}

\providecommand{\nichttitle}{Formelsammlung}
\providecommand{\nichtsubtitle}{}

\renewcommand{\nichttitle}{Numerische Methoden}
\renewcommand{\nichtsubtitle}{Formelsammlung}


\begin{document}

\begin{center}
{\small\bfseries\color{accent}\nichttitle}
\end{center}
\vspace{-6pt}

\scriptsize
\begin{multicols}{2}

% ── BODY ──────────────────────────────────────────────────────

\section{Numerische Methoden — Spickzettel}


\subsection{1. Grundlagen}

\textbf{Numerische Methoden}: näherungsweise Lösung mathematisch nicht/unpraktikabel analytisch lösbarer Probleme.

\begin{itemize}
  \item f(·)=exakt, f̃(·)=Näherung; x=exakt, x̃=gestört
  \item Nötig bei: nichtlinearen Gln., transzendenten Fkt., großen LGS
\end{itemize}

\textbf{Transzendente Funktion}: keine endliche Kombi algebraischer Operationen → keine geschlossene Lösung. Bsp. σ(x)=1/(1+e⁻ˣ)

\textbf{Smooth}: Ableitungen aller Ordnungen existieren.

\textbf{Moore's Law}: Rechenleistung verdoppelt sich \textasciitilde{}alle 2 Jahre.

\textbf{[a,b]}: geschlossenes Intervall (inkl. Endpunkte).



\subsection{2. Diskretisierung}

\textbf{Diskretisierung}: kontinuierliche Domäne → endliche Schritte/Gitter.

\begin{itemize}
  \item Kontinuierlich f(x), x∈[a,b] → Diskret f(xᵢ), xᵢ=a+ih, i=0,…,N
  \item Schrittweite \textbf{h = (b−a)/N}
\end{itemize}


\subsection{3. Floating-Point Arithmetik}


\subsubsection{3.1 Fehlertypen}

{\small
\begin{adjustbox}{max width=\linewidth}
{\scriptsize
\begin{tabular}{ll}
\toprule
\textbf{Fehler} & \textbf{Ursache} \\
\midrule
Modeling Error & Modell vereinfacht Realität \\
Truncation/Approximation & endliche Approx. unendlicher Prozesse \\
Rounding Error & endliche Stellenzahl \\
Overflow & Magnitude zu groß → ±∞ \\
Underflow & Magnitude nahe 0 → 0 \\
\bottomrule
\end{tabular}
}
\end{adjustbox}
}



\subsubsection{3.2 Floating-Point Darstellung}
\textbf{x = (−1)ˢ · m · 2ᵉ} (IEEE 754); s=Vorzeichen, m=Mantisse, e=Exponent.



{\small
\begin{adjustbox}{max width=\linewidth}
{\scriptsize
\begin{tabular}{lllll}
\toprule
\textbf{Format} & \textbf{Bits} & \textbf{Sign} & \textbf{Exponent} & \textbf{Mantisse} \\
\midrule
HALF & 16 & 1 & 6 & 9 \\
SINGLE & 32 & 1 & 8 & 23 \\
DOUBLE & 64 & 1 & 11 & 52 \\
\bottomrule
\end{tabular}
}
\end{adjustbox}
}


\begin{itemize}
  \item IEEE 754 Single: Bias=127, e=E−127; E∈0…255
  \item Größte Single ≈ 10³⁸ (log₁₀(2¹²⁸)≈38.5)
  \item Mantisse t Bits → 2ᵗ Werte zwischen zwei Zweierpotenzen
  \item 2-Bit Bsp: 0.00=0, 0.01=0.25, 0.10=0.5, 0.11=0.75
\end{itemize}

\textbf{Machine Epsilon ε\_mach}: kleinste pos. Zahl mit 1+ε\_mach≠1.

\begin{itemize}
  \item \textbf{ε\_mach = 2⁻ᵗ}
  \item Single ≈ 1.19×10⁻⁷; Double ≈ 2.22×10⁻¹⁶
  \item Absolute Präzision: ε\_mach·|x| (große x → größere Lücken)
\end{itemize}

\textbf{Fixed-Point}: feste Bits Integer-/Fraktionsteil → konstante Präzision.

\begin{itemize}
  \item Skalierung 10⁶: 1.234567 → 1234567; max. Rundungsfehler 0.5×10⁻⁶
\end{itemize}


\subsubsection{3.3 Loss of Significance}
\textbf{Catastrophic Cancellation}: Subtraktion fast gleicher Zahlen löscht führende Stellen.

\begin{itemize}
  \item 8 sig. Digits: a=12345678.5, b=12345678.0 → ã−b̃=1 statt 0.5
  \item f(ε)=1/(1+ε−1) divergiert bei kleinem ε
\end{itemize}

\textbf{Pseudo-Accuracy}: künstliches Detail suggeriert ungerechtfertigte Genauigkeit.



\subsection{4. Fehleranalyse}


{\small
\begin{adjustbox}{max width=\linewidth}
{\scriptsize
\begin{tabular}{lll}
\toprule
\textbf{Begriff} & \textbf{Formel} & \textbf{Bedeutung} \\
\midrule
\textbf{Conditioning} κ & κ=\textbackslash{} & f(x)−f(x̃)\textbackslash{} \\
\textbf{Stability} & \textbackslash{} & f̃(x̃)−f̃(x)\textbackslash{} \\
\textbf{Consistency} & \textbackslash{} & f̃(x)−f(x)\textbackslash{} \\
\textbf{Convergence} & \textbackslash{} & f̃(x̃)−f(x)\textbackslash{} \\
\bottomrule
\end{tabular}
}
\end{adjustbox}
}


\begin{itemize}
  \item κ unabhängig von Methode
  \item Konvergenz meist →0; Nicht-Null = systematischer Bias
  \item Stability = Conditioning des Systems für h→0
\end{itemize}


\subsection{5. Brute-Force / Grid Search}

\textbf{Grid Search}: f an allen Gitterpunkten auswerten, Bestes wählen.

\begin{itemize}
  \item Komplexität \textbf{O(Nᵈ)}
  \item Extrema auf [a,b]: kritische Punkte ODER Endpunkte prüfen
\end{itemize}


\subsubsection{5.1 min sin(x) auf [−3,1]}
\begin{itemize}
  \item Analytisch: cos(x)=0 ⟹ x*=−π/2≈−1.5708, sin(−π/2)=−1
  \item Grid h=1: sin(−3)≈−0.1411, sin(−2)≈−0.9093, sin(−1)≈−0.8415, sin(0)=0, sin(1)≈0.8415
  \item Num. Min: −0.9093 bei x=−2; Approx.fehler 0.0907
\end{itemize}


\subsubsection{5.2 LGS-Beispiel}
[2,1;5,1]·[w,b]ᵀ=[11,13]ᵀ → w=2/3, b=29/3

\begin{itemize}
  \item Grid w,b∈\{0,2.5,5,7.5,10\}; error₁=|2w+b−11|, error₂=|5w+b−13|
  \item Min bei w=0, b=10, akkum. Fehler=4
  \item Substitution skaliert nicht für große Systeme
\end{itemize}


\subsubsection{5.3 Conditioning/Stability min sin(x), x̃=x±0.25}
\begin{itemize}
  \item κ(−1,+0.25)≈0.1599; κ(−1,−0.25)≈0.1074 (gut); κ(0,±0.25)≈0.2474 (schlecht)
  \item Stability h=1: s≈0; h=0.2 bei x=−1: s≈0.4964
  \item Consistency: c₁≥0.15853; c₀.₂≥0.05102 (besser)
  \item Beide h gleich konvergent; h=0.2 nutzt nichts wegen Stabilität
\end{itemize}


\subsection{6. Taylor-Expansion}

\[
f(x)=\sum_{k=0}^{\infty}\frac{f^{(k)}(x_n)}{k!}(x-x_n)^k
\]



{\small
\begin{adjustbox}{max width=\linewidth}
{\scriptsize
\begin{tabular}{lll}
\toprule
\textbf{Ordnung} & \textbf{P(x)} & \textbf{Bedingung} \\
\midrule
0. & f(xₙ) & P(xₙ)=f(xₙ) \\
1. & f(xₙ)+f'(xₙ)(x−xₙ) & P'(xₙ)=f'(xₙ) \\
2. & + f''(xₙ)/2·(x−xₙ)² & P''(xₙ)=f''(xₙ) \\
\bottomrule
\end{tabular}
}
\end{adjustbox}
}


\begin{itemize}
  \item /k! kompensiert Mehrfachableitung
  \item \textbf{Lokale} Approx., verankert an xₙ
\end{itemize}


\subsection{7. Newton-Raphson}

\textbf{Ziel}: f(x*)=0 via lokaler Steigung.

\textbf{Lin. Approx.}: f(x)≈f(xₙ)+f'(xₙ)(x−xₙ)


\[
x_{n+1}=x_n-\frac{f(x_n)}{f'(x_n)}
\]


\begin{quote}
Merke: Linear reicht, weil wir f(x*)=0 suchen, nicht f(x) approximieren.
\end{quote}


\subsubsection{7.1 Algorithmus 1}
\textbf{Require}: x₀, f, f', ε

\begin{enumerate}
  \item x ← x₀
  \item while |f(x)|>ε:
  \item   d ← f'(x)
  \item   if |d|<ε\_machine: error "Derivative too small"
  \item   x ← x − f(x)/d
  \item return x
\end{enumerate}


\subsubsection{7.2 Quadratische Konvergenz}
\begin{itemize}
  \item eₙ=xₙ−x\textit{; f(xₙ)=f'(x})·eₙ+f''(ξ)/2·eₙ²
  \item $e_{n+1}\approx-\frac{f''(\xi)}{2f'(x^*)}e_n^2$
  \item \textbf{|eₙ₊₁|≤C·|eₙ|²}; Fehler quadriert: 10⁻¹→10⁻²→10⁻³→10⁻⁶→10⁻¹²
\end{itemize}


\subsubsection{7.3 √2, x₀=2}

{\small
\begin{adjustbox}{max width=\linewidth}
{\scriptsize
\begin{tabular}{lllll}
\toprule
\textbf{n} & \textbf{xₙ} & {\color{white}\textbf{\textbackslash{}}} & {\color{white}\textbf{eₙ\textbackslash{}}} & \textbf{} \\
\midrule
0 & 2.000000000 & 5.86×10⁻¹ &  &  \\
1 & 1.500000000 & 8.58×10⁻² &  &  \\
2 & 1.416666667 & 2.45×10⁻³ &  &  \\
3 & 1.414215686 & 2.13×10⁻⁶ &  &  \\
4 & 1.414213562 & 4.5×10⁻¹² &  &  \\
\bottomrule
\end{tabular}
}
\end{adjustbox}
}



\subsubsection{7.4 sin(x)=0, x₀=3.5}
\begin{itemize}
  \item x₁=3.5−(−0.351)/(−0.936)=3.125
  \item x₂≈3.14159; nach 2 Iter. Fehler<10⁻⁵
  \item Vgl. Grid h=0.3: 6 Auswertungen, Fehler≈0.042
\end{itemize}


\subsubsection{7.5 f(x)=x³−x, x₀=1.4 (f'=3x²−1)}
\begin{itemize}
  \item x₁=1.4−1.344/4.88≈1.127
  \item x₂≈0.976; x₃≈1.014
\end{itemize}


\subsubsection{7.6 Failure Modes}
\begin{itemize}
  \item \textbf{Zero derivative}: f'(xₙ)=0 → horizontale Tangente
  \item \textbf{Oszillation}: symmetrische Fkt. → Zyklus
  \item \textbf{Ambiguous starting point}: x₀ entscheidet Wurzel
  \item \textbf{Lokales Optimum}: nahes ≠ globales
  \item \textbf{Divergenz}: bei f' nahe 0
\end{itemize}


\subsection{8. Secant Methode}

\textbf{Motivation}: f'(x) nicht verfügbar (teuer/Black-Box/nicht diff.bar).

\textbf{Approx. f'(xₙ)} ≈ (f(xₙ)−f(xₙ₋₁))/(xₙ−xₙ₋₁)


\[
x_{n+1}=x_n-f(x_n)\cdot\frac{x_n-x_{n-1}}{f(x_n)-f(x_{n-1})}
\]


\begin{itemize}
  \item \textbf{2 Startwerte} x₀, x₁
  \item Geometrisch: Sekantensteigung statt Tangente
\end{itemize}


\subsubsection{8.1 Algorithmus 2}
\textbf{Require}: x₀, x₁, f, ε

\begin{enumerate}
  \item while |f(x₁)|>ε:
  \item   s ← (f(x₁)−f(x₀))/(x₁−x₀)
  \item   x\_old ← x₁
  \item   x₁ ← x₁ − f(x₁)/s
  \item   x₀ ← x\_old
  \item return x₁
\end{enumerate}


\subsubsection{8.2 f(x)=x³−x, x₀=1.2, x₁=1.4}
\begin{itemize}
  \item x₂≈1.4−0.744·(0.2/0.616)≈1.159
  \item x₃≈1.011
\end{itemize}


\subsection{9. Vergleich}


{\small
\begin{adjustbox}{max width=\linewidth}
{\scriptsize
\begin{tabular}{lllll}
\toprule
\textbf{Methode} & \textbf{Info} & \textbf{Startwerte} & \textbf{Konvergenz} & \textbf{Komplexität} \\
\midrule
Grid Search & f(x) & Intervall & linear/blind & O(Nᵈ) \\
Newton-Raphson & f(x), f'(x) & 1 (x₀) & quadratisch & wenige Iter. \\
Secant & f(x) & 2 (x₀,x₁) & superlinear & kein f' nötig \\
\bottomrule
\end{tabular}
}
\end{adjustbox}
}



\subsection{10. ML/DL Bezug}

\begin{itemize}
  \item Newton 1. Ordnung = \textbf{Gradient Descent}
  \item \textbf{Backpropagation}: verteilt Gradient durch Netz-Schichten
  \item Lokale Methoden → lokale Lösungen; bei f'=0 oszillieren/divergieren
\end{itemize}






% ──────────────────────────────────────────────────────────────

\end{multicols}
\end{document}
